\subsection{健身、骑行与性功能：运动的"双刃剑"}

规律运动总体有益于性功能，但\textbf{运动的类型、强度与姿势}也会带来另一面——某些练法反而可能损害勃起与性欲。本节聚焦最常被误解的几种情况。

\vspace{0.5em}
\noindent\textbf{1. 力量训练与睾酮：别指望"练出激素"}

\begin{itemize}
	\item 单次大强度训练后睾酮可有\textbf{短暂升高}，但这是应激反应，会很快回落；长期看，规律力量训练更多是通过改善体成分、胰岛素敏感性与自信来间接提升性功能，\textbf{而非把睾酮"练高"}。
	\item \textbf{过度训练反而有害}：长期高频、大重量、睡眠不足且热量摄入不足的"过度训练综合征"，会抑制下丘脑--垂体--性腺轴，导致\textbf{睾酮下降、性欲减退、晨勃减少}。若训练后持续疲劳、情绪低落、性欲明显下降，应减量并保证睡眠与营养。
	\item \textbf{合成代谢类固醇}：外源雄激素会抑制自身睾酮分泌，长期滥用可致睾丸萎缩、性欲紊乱与不育（详见"药物与性功能速查"章）。
\end{itemize}

\noindent\textbf{2. 骑行：会阴压迫是真实风险}

\begin{itemize}
	\item 长时间骑行时身体重量压在狭窄车座上，可压迫\textbf{会阴部血管与阴部神经}，导致会阴麻木、勃起功能下降，即所谓"骑行相关性功能障碍"。
	\item \textbf{降低风险的做法}：选择带中央凹槽或镂空的\textbf{减压车座}，适当抬高车把以减少前倾、把体重分散到坐骨；穿\textbf{带护垫的骑行裤}（不穿内裤以减少摩擦）；每骑行 20--30 分钟\textbf{起身放松}一次；控制单次连续骑行时长。
	\item \textbf{若出现会阴麻木、勃起变差}，应减少骑行、更换装备，并就医评估；多数在调整后改善。
\end{itemize}

\noindent\textbf{3. 运动性热暴露：桑拿、热水浴与久坐}

\begin{itemize}
	\item 精子对温度敏感。频繁\textbf{桑拿、热水坐浴、长时间泡热水澡}会升高阴囊温度、抑制精子生成；备孕男性应适度。
	\item 同理，久坐与笔记本长时间置于腿上也是常见的"热暴露"来源（详见"环境因素与性健康"）。
\end{itemize}

\noindent\textbf{4. 对性功能最"友好"的练法}

\begin{itemize}
	\item \textbf{有氧运动}：每周 150 分钟中等强度（快走、慢跑、游泳、骑行需注意座垫），改善血管内皮功能，是勃起功能的"基础工程"。
	\item \textbf{盆底肌训练（凯格尔）}：男女均受益——男性改善勃起硬度与射精控制，女性改善性唤起与高潮（详见"盆底健康与性功能"章）。
	\item \textbf{力量训练}：每周 2 次全身训练，改善体成分与自信，间接提升性欲与耐力。
	\item \textbf{瑜伽、拉伸与放松}：减轻压力、改善身体觉察，对性欲与性满意度有帮助。
\end{itemize}

\begin{tcolorbox}[colback=pink!5!white,colframe=red!50!black,title=贴心小叮咛]
运动是性功能的"长期投资"，不是"立竿见影的春药"。\textbf{练对了加分，练过了减分，骑车不换座垫可能直接扣分}。若出现会阴麻木、晨勃减少或性欲明显下降，先检查训练量、睡眠与装备，再决定是否就医。
\end{tcolorbox}
